\documentclass[fleqn, usenatbib, useAMS]{mnras}
\usepackage{graphicx}
\usepackage{amsmath,paralist,xcolor,xspace,amssymb}
\usepackage{times}
\usepackage[super]{nth}

\newcommand{\swift}{{\sc swift}\xspace}
\newcommand{\gadget}{{\sc gadget}\xspace}
\newcommand{\dd}[2]{\frac{\partial #1}{\partial #2}}
\renewcommand{\vec}[1]{{\mathbf{#1}}}
\newcommand{\Wij}{\overline{\nabla_xW_{ij}}}
\newcommand{\tbd}{\textcolor{red}{TO BE DONE}}

\newcommand{\MinimalSPH}{\textsc{minimal-sph}\xspace}
\newcommand{\GadgetSPH}{\textsc{gadget-sph}\xspace}
\newcommand{\PESPH}{\textsc{pe-sph}\xspace}
%\setlength{\mathindent}{0pt}

%opening
\title{ULDM implementation in \swift}
\author{Yves Revaz}

\begin{document}

\maketitle
\section{Kernel functions}


\subsection{Kernel Derivatives}

The first derivatives of the kernel function have relatively simple
expressions:

\begin{align}
 \vec\nabla_x W(\vec{x},h) &= \frac{1}{h^4}f'\left(\frac{|\vec{x}|}{h}\right) \frac{\vec{x}}{|\vec{x}|}, \\
 \frac{\partial W(\vec{x},h)}{\partial h} &=- \frac{1}{h^4}\left[3f\left(\frac{|\vec{x}|}{h}\right) + 
\frac{|\vec{x}|}{h}f'\left(\frac{|\vec{x}|}{h}\right)\right].
\end{align}
Note that for all the kernels listed above, $f'(0) = 0$. 


The second derivatives of the kernel write:

\begin{align}
 \vec\nabla_x W(\vec{x},h) &= \frac{1}{h^4}f'\left(\frac{|\vec{x}|}{h}\right) \frac{\vec{x}}{|\vec{x}|}, \\
 \frac{\partial W(\vec{x},h)}{\partial h} &=- \frac{1}{h^4}\left[3f\left(\frac{|\vec{x}|}{h}\right) + 
\frac{|\vec{x}|}{h}f'\left(\frac{|\vec{x}|}{h}\right)\right].
\end{align}







\section{Derivation of the Equation of Motion}\label{sec:derivation}
The implementation follows the ones given by \citet{nori2024}.

For a set of particles $i$ with positions $\vec{x}_i$, masses $m_i$,
and smoothing length $h_i$, we compute the density for each particle as:
%
\begin{equation}
  \rho_i \equiv \rho(\vec{x}_i) = \sum_j m_j W(\vec{x}_{ij}, h_i).
  \label{eq:uldm:rho}
\end{equation}
and the derivative of its density with respect to $h$:

\noindent
Several possibilities exists to write the density gradient.
The standard way writes:
\begin{equation}
  \label{eq:uldm:grad_rho}
  \vec{\nabla} \rho_i =  \sum_j \frac{m_j}{\rho_j} \left( \rho_j - \rho_i \right) \vec{\nabla} W(\vec{x}_{ij}, h_i),
\end{equation}
which includes the substraction of the first order term. 
Another choice is \citep{nori2024}:
%
\begin{equation}
  \label{eq:uldm:grad_rho}
  \vec{\nabla} \rho_i =  \sum_j \frac{m_j}{\sqrt{\rho_i \rho_j}} \left( \rho_j - \rho_i \right) \vec{\nabla} W(\vec{x}_{ij}, h_i).
\end{equation}
This latter provides slightly better results on the 1D density front.

\noindent
As for the gradient density, several choices to write the Laplacian exits.
The simplest formulation  \citep[][Eq.~86]{price2012}:
%
\begin{equation}
  \nabla^2 \rho_i =  \sum_j \frac{m_j}{\rho_j} \rho_j \vec{\nabla}^2 W(\vec{x}_{ij}, h_i) =  \sum_j m_j  \vec{\nabla}^2 W(\vec{x}_{ij}, h_i),
\end{equation}
or removing the first order term \citep[][Eq.~90]{price2012}:
%
\begin{equation}
  \nabla^2 \rho_i =  \sum_j \frac{m_j}{\rho_j} \left( \rho_j - \rho_i \right)  \vec{\nabla}^2 W(\vec{x}_{ij}, h_i).
\end{equation}
Note that those two formulations totally fails the 1D density front test.
An alternative is to compute the Laplacian from the gradient (this works only in 1D I guess).
%
\begin{equation}
  \nabla^2 \rho_i =  \sum_j \frac{m_j}{\sqrt{\rho_i \rho_j}} \left( \vec{\nabla}\rho_j - \vec{\nabla}\rho_i \right) \vec{\nabla} W(\vec{x}_{ij}, h_i).
\end{equation}
This formulation is working.
Another formulation is \citep[][Eq.~18]{nori2018} {\bf (check with Price eq. 91)}:
%
\begin{equation}
  \nabla^2 \rho_i =  -2 \sum_j m_j \frac{ \rho_j - \rho_i }{\rho_j} \frac{ \vec{r_{ij}} \vec{\nabla} W(\vec{x}_{ij}, h_i)}{r^2_{ij}}.
\end{equation}





The Laplacian of the density write:
%
\begin{equation}
  \label{eq:uldm:laplacian_rho}
  \nabla^2 \rho_i =  \sum_j m_j \vec{\nabla}^2 W(\vec{x}_{ij}, h_i) \frac{\rho_j - \rho_i}{\sqrt{\rho_i \rho_j}} - \frac{|\vec{\nabla}\rho_i|^2}{\rho_i}.
\end{equation}

\noindent
The quantum potential of the particle $i$ is :
%
\begin{equation}
  \label{eq:uldm:grad_Q}
  Q_i =  \frac{\hbar}{2 m^2_\chi }  \sum_j \frac{m_j}{\rho_j} W(\vec{x}_{ij}, h_i)
  \left( \frac{\nabla^2\rho_j}{2\rho_j} -\frac{|\vec{\nabla}\rho_j|^2}{4\rho^2_j} \right)
\end{equation}


The acceleration of the particle $i$ is equal to the gradient of the quantum potential:

\begin{equation}
  \label{eq:uldm:grad_Q}
  \vec{\nabla{Q_i}} =  \frac{\hbar}{2 m^2_\chi }  \sum_j \frac{m_j}{\rho_j}\vec{\nabla} W(\vec{x}_{ij}, h_i)
  \left( \frac{\nabla^2\rho_j}{2\rho_j} -\frac{|\vec{\nabla}\rho_j|^2}{4\rho^2_j} \right)
\end{equation}





\section{Implementation}\label{sec:implementation}
In the following, the function $f(u)$, $f'(u)$ and $f''(u)$ are respectively
the kernel and its first and second derivative (excluding $h$).
They are provided by the the \texttt{kernel\_deval} and \texttt{kernel\_d2eval} functions.

The density is computed in the first loop, as:

\begin{equation}
  \rho_i = \sum_j m_j W(\vec{x}_{ij}, h_i) = \frac{1}{h_i^3} \sum_j  m_j f(u_i), 
  \label{eq:uldm:implementation:rho}
\end{equation}

with $u_i = r/h_i $ and $r=|\vec{x}_i-\vec{x}_j|$.
Note that the $\frac{1}{h_i^3}$ factor is computed at the end of the loop.


The gradient of the density is computed in the second loop as:

\begin{eqnarray}
  \label{eq:uldm:implementation:grad_rho}
  \vec{\nabla} \rho_i & = &  \sum_j m_j \vec{\nabla} W(\vec{x}_{ij}, h_i) \frac{\rho_j - \rho_i}{\sqrt{\rho_i \rho_j}} \\
                      & = & \frac{1}{h_i^4}  \sum_j m_j f'(u_i) \frac{\vec{x}}{|\vec{x}|} \frac{\rho_j - \rho_i}{\sqrt{\rho_i \rho_j}}
\end{eqnarray}

Note that the $\frac{1}{h_i^4}$ factor is computed at the end of the loop.


The Laplacian of the density is computed in the second loop, as:

\begin{eqnarray}
  \label{eq:uldm:implementation:laplacian_rho}
  \nabla^2 \rho_i
  & = &  \sum_j m_j \vec{\nabla}^2 W(\vec{x}_{ij}, h_i) \frac{\rho_j - \rho_i}{\sqrt{\rho_i \rho_j}} - \frac{|\vec{\nabla}\rho_i|^2}{\rho_i}\\
  & = & \frac{1}{h_i^5}  \sum_j m_j \left[ f''(u_i) + \frac{2}{u_i} f'(u_j) \right]  \frac{\rho_j - \rho_i}{\sqrt{\rho_i \rho_j}} \\
  && - \frac{|\vec{\nabla}\rho_i|^2}{\rho_i}
\end{eqnarray}

Note that the  $\frac{1}{h_i^5}$ factor as well as the last ratio that involve only $\rho_i$ are computed at the end of the loop.


The gradient of the quantum pressure is computed in the third loop, as:

\begin{eqnarray}
  \label{eq:uldm:implementation:grad_Q}
  \vec{\nabla{Q_i}}
  & = &  \frac{\hbar}{2 m^2_\chi }  \sum_j \frac{m_j}{\rho_j}\vec{\nabla} W(\vec{x}_{ij}, h_i)
  \left( \frac{\nabla^2\rho_j}{2\rho_j} -\frac{|\vec{\nabla}\rho_j|^2}{4\rho^2_j} \right)\\
  & = &  \frac{\hbar}{2 m^2_\chi } \frac{1}{h_i^4}
         \sum_j \frac{m_j}{\rho_j} f'(u_i) \frac{\vec{x}}{|\vec{x}|}
         \left( \frac{\nabla^2\rho_j}{2\rho_j} -\frac{|\vec{\nabla}\rho_j|^2}{4\rho^2_j} \right)
\end{eqnarray}

Note that the  $\frac{1}{h_i^4}$ factor as well as the constant are computed at the end of the loop.










\bibliographystyle{mnras}
\bibliography{./bibliography}

\end{document}
